% SOURCES: README.md (Stack); CLAUDE.md (Tech Stack); docs/0.2.0v/plan.md
%          (§9 Library-to-Feature Map — the JUSTIFICATION source);
%          docs/1.0.0v/supervisor-report.md

\chapter{Technologies Used}
\label{ch:tech}

\section{Overview}
The application is a single TypeScript codebase spanning three runtime contexts:
an Electron \emph{main} process (privileged, native OS access), a \emph{preload}
security bridge, and a React \emph{renderer} (the editor UI). It is built and
packaged with Vite and electron-builder. Every technology was chosen against a
specific requirement rather than by default; Table~\ref{tab:tech} lists the
principal choices and the reason for each, and the sections that follow expand on
the choices that best show the engineering rationale.

\begin{table}[H]
  \centering
  \caption{Principal technologies and justification.}
  \label{tab:tech}
  \small
  \begin{tabularx}{\linewidth}{@{}l l X@{}}
    \toprule
    \textbf{Technology} & \textbf{Role} & \textbf{Justification} \\
    \midrule
    Electron \cite{electron} & Desktop shell & Cross-platform desktop with native file access; local-first, no server required. \\
    React 18 \cite{react} & UI & Declarative rendering of the document tree; mature ecosystem and tooling. \\
    TypeScript 5 \cite{typescript} & Language & Strict typing across three processes; the cross-owner contracts are enforced at compile time. \\
    Zustand + immer \cite{zustand,immer} & State + history & Minimal store with no boilerplate; \texttt{immer} patches yield undo/redo almost for free. \\
    Zod \cite{zod} & Validation & Schema validation at the IPC and file-load boundaries; types are derived via \texttt{z.infer} to keep schema and type in lockstep. \\
    axe-core \cite{axecore} & A11y gate & Industry-standard accessibility engine; blocks export on any critical or serious issue. \\
    prettier \cite{prettier} & Output format & Formats generated HTML/CSS so it reads as hand-written before minification. \\
    Lightning CSS \cite{lightningcss} & CSS minify & Fast, correct CSS transform and minification in the export pipeline. \\
    sharp \cite{sharp} & Images & Generates WebP and \texttt{srcset} variants in the main process. \\
    Vite \cite{vite} & Build / dev & Fast hot-module-reload dev server and production bundling. \\
    MCP SDK \cite{mcp-spec} & Headless API & Exposes the pipeline as tools to agents and clients over \texttt{stdio}. \\
    \bottomrule
  \end{tabularx}
\end{table}

Beyond the principals, the editor uses \texttt{@dnd-kit} for canvas drag-and-drop
and \texttt{react-arborist} for the layers tree; Radix UI primitives for accessible
tabs, dialogs, and popovers; \texttt{react-colorful} and \texttt{chroma-js} for the
colour picker and contrast checks; and \texttt{jszip} plus
\texttt{html-minifier-terser} in the export pipeline. Each of these maps to a
specific Tier-1 feature in the execution plan's library-to-feature justification.

\section{Selected Justifications in Depth}

\subsection*{Why a single Document Model backed by immer}
Making one immutable \texttt{Document} the source of truth means that undo/redo,
persistence, and generation all reduce to functions of one value. \texttt{immer}
is what makes this ergonomic: each operation is written as if it were mutating a
draft, while \texttt{produceWithPatches} captures the forward and inverse patches
that the history store replays for undo and redo. The operations therefore never
need to know anything about history --- reversibility is a property of the patch
round-trip, not of the operation code.

\subsection*{Why an accessibility hard gate rather than a linter}
Accessibility is treated as a build \emph{error}, not a warning. The
\texttt{axe-core} engine runs on the generated HTML inside a \texttt{jsdom}
document, and any critical or serious violation aborts the export. This guarantees
a compliance floor that no manual review process could match consistently, and ---
crucially --- it applies to \emph{every} export path, including the headless MCP
server, so the guarantee cannot be bypassed by scripting the tool.

\subsection*{Why prettier plus stable identifiers gives deterministic output}
A recurring requirement is that identical input must yield byte-identical output.
Two decisions make this hold: element identifiers are generated once and are stable
across save/load (they are never regenerated), and the generator emits no
timestamps or random values. Running the emitted HTML and CSS through
\texttt{prettier} then produces a canonical, human-readable form, so two exports of
the same document are not merely equivalent but textually identical --- which makes
the output diff-able and the pipeline testable by exact comparison.
